\documentclass[a4paper,10pt]{article}
\usepackage[francais]{babel}
\usepackage{url}

\input{letter-icube-uds.sty}

\enCouleur

\begin{document}

\name{Stéphane \textsc{Genaud}, Philippe \textsc{Helluy} }
%\fonction{Professeur ENSIIE, mise à disposition Université de Strasbourg}
\telephoneillkirch{+33 (0)3688 54542}
\email{genaud@unistra.fr, helluy@unistra.fr}

\begin{letter}[Recommandation de Clément Flint]{à}{Ecole Doctorale MSII}

Chers collègues, 

  Nous souhaitons co-encadrer  une thèse ayant pour  objectif l'amélioration des
  méthodes de  compression des données  dans le  domaine du HPC,  en particulier
  pour la simulation numérique de phénomènes physiques.

Pour ce  travail, Clément Flint  s'est porté candidat  et nous avons  soutenu au
plus haut  point sa  candidature dans  les campagnes  de financement  de l'Ecole
Doctorale MSII et du LABEX  IRMIA. Sa candidature a été retenue à la fois par
l'Ecole Doctorale et le LABEX, et il a choisi de prendre le financement LABEX.

Clément est  un étudiant  qui a  obtenu l'an dernier  le Master  en informatique
(filière Cursus  Master en Ingénierie  IIRVIJ, en  se classant major  toutes les
cinq années de son cursus.  Il a  effectué son stage orienté recherche de Master
dans notre équipe  Icube-ICPS sous la supervision de  Bérenger Bramas, principal
encadrant dans la thèse à venir.

Le  travail  de  thèse  sera  pluri-discplinaire, à  la  fois  en  mathématiques
appliquées à des  modélisations de phénomènes physique  et informatique. Clément
sera  encadré en  co-direction  de  thèse par  Philippe  Helluy  pour la  partie
applicative et  par Stéphane  Genaud et Bérenger  Bramas (co-encadrant)  pour la
partie informatique. Clément serait hébergé dans  les bureaux de l'équipe ICPS à
ICube,  mais les  interractions avec  l'équipe \textit{Modélisation  et contrôle}
de l'IRMA sont fréquentes.


Nous sommes convaincus que Clément Flint est un excellent candidat pour ce
travail, et nous soutenons encore une fois sans réserve sa candidature pour une
inscription en thèse.

\begin{flushright}
\includegraphics[width=.25\textwidth]{signgenaud.jpg} 
\includegraphics[width=.25\textwidth]{sign_helluy.jpg} 
\end{flushright}

\end{letter}
\newpage
\noindent
Avis du directeur de l'unité de recherche, Michel \textsc{de Mathelin}:\\[3mm]
%\vspace{3mm}
\framebox{
\begin{minipage}{.9\linewidth}
\hspace{.9\linewidth}
\vspace{4.3cm}

\end{minipage}
}

%\newpage
%\bibliographystyle{plain}
%\bibliography{arno}


\end{document}
